% Source: https://artofproblemsolving.com/community/c3h2443434_number_theory_problems
% Stable identity: aops-topic-2443434-q03
Sort the $n$ numbers $1,2,\dots,n$ according to their greatest common divisor with $n$. For a positive divisor $g$ of $n$, let $A_g$ be the set of $k\in\{1,\dots,n\}$ with $\gcd(k,n)=g$. Every $k$ lies in exactly one $A_g$, so
\[\sum_{g\mid n}|A_g|=n.\]

\textbf{Counting $A_g$.} If $\gcd(k,n)=g$, then $k=gj$ with $1\le j\le\dfrac ng$, and $\gcd(gj,n)=g\cdot\gcd\left(j,\dfrac ng\right)$, so $\gcd(k,n)=g$ exactly when $\gcd\left(j,\dfrac ng\right)=1$. Hence $k\mapsto\dfrac kg$ is a bijection from $A_g$ to the integers $j\le\dfrac ng$ coprime to $\dfrac ng$, and $|A_g|=\phi\left(\dfrac ng\right)$.

Therefore $n=\sum_{g\mid n}\phi\left(\dfrac ng\right)$. As $g$ runs over the divisors of $n$, so does $d=\dfrac ng$, and the sum is $\sum_{d\mid n}\phi(d)$. $\blacksquare$

(Example: for $n=12$, $\phi(1)+\phi(2)+\phi(3)+\phi(4)+\phi(6)+\phi(12)=1+1+2+2+2+4=12$.)
